\chapter{Inference}
\label{sec:inference}

\section{Shape}
\label{sec:inf-shape}

The analysis runs per module, after type checking, at the point where the effect inference runs
today: every type is solved, and nothing has yet been published. It cannot run later, because
beni's backend sees no types. For each module it makes four passes:

\begin{enumerate}
\item \emph{Local flow} (\S\ref{sec:inf-flow}): one forward walk of each declaration's body in
evaluation order, computing $A(e)$ for every tracked expression. Calls use the callees' summaries;
a callee in the same strongly connected component uses the current round's summary.
\item \emph{Liveness} (\S\ref{sec:inf-liveness}): one backward walk per declaration, computing
$\live(h, p)$ for the holders the forward walk named.
\item \emph{Check and summarise} (\S\ref{sec:inf-summarise}): the four side conditions at every
demand, at every call of a once closure and at every creation of one, and the declaration's summary
read off the two walks.
\item \emph{The group fixpoint} (\S\ref{sec:inf-recursion}): passes 1--3 repeated for a recursive
group until no summary changes, followed by a validation round.
\end{enumerate}

Function-typed values carry \emph{ownership classes} (\S\ref{sec:inf-classes}), solved with
directional edges in the same pass. Imported summaries come from interfaces
(\S\ref{sec:inf-modular}). The pass runs on every build, as the effect pass does, and, unlike the
write-set analysis, it produces diagnostics.

\section{Local flow}
\label{sec:inf-flow}

beni's intermediate representation is already a sequence of instructions in evaluation order, with
\code{let}-style binding and \code{case} branching; the walk reads it as the normal form of
\S\ref{sec:model-unique}, in which every argument and every field of a structure under construction
is a variable, used at the instruction that consumes it. The forward walk keeps $\Sigma$ and applies
the transfer rules of Table~\ref{tab:transfer}; a \code{case} forks $\Sigma$ per arm and joins the
results by union. The only loop inside a body is a self tail call, which the walk treats as a
recursive call of the same declaration: the parameters' cells on re-entry are the union of the entry
cells and the jump's arguments, an inner fixpoint over the body that converges in at most as many
rounds as there are distinct abstract cells in the body.

Directionality and flow sensitivity both matter, as one example shows:
\begin{Verbatim}
xs = [ 1, 2, 3 ]          -- A(xs) = {site₁}
ys = xs                   -- A(ys) = {site₁}: a copy of the set, not a merged class
zs = List.set ys 0 9      -- demand on site₁; holders xs and ys; is xs dead after?
List.length zs            -- yes: xs is never used again. Accepted.
\end{Verbatim}
A unification-based model would put \code{xs}, \code{ys} and \code{zs} in one class, and any later use
of any of them would taint all three. Here \code{zs} is a new holder of the \emph{same} cell---ownership
is transferred---and \code{xs} is dead.

The cost of the forward walk is linear in the number of instructions times the width of the
abstract values (\S\ref{sec:inf-complexity}). Cell sets are usually singletons. Abstract values are
hash-consed, as the write-set analysis's are, so a chain of conditionals does not duplicate them.

\section{Liveness}
\label{sec:inf-liveness}

Liveness is a standard backward pass over the same instructions, with one bit per (variable, path)
and the uses of \S\ref{sec:model-liveness}. Branches join by union: a holder live in either arm is
live before the \code{case}. A captured path is live \emph{wherever the closure value is live}---at
every call, store, return or pass of the closure, or of any structure that holds it, which the walk
sees as ordinary uses of the holder's paths through $\env$---and \emph{from its creation onward,
forever}, when the closure escapes (is passed to a position whose $\mathit{use}$ contains
$\mathsf{E}$ at $\env$, or is stored in an escaped place). So
\begin{Verbatim}
g = λ_ → List.length xs
ys = List.push xs 1
{ m | f = g }
\end{Verbatim}
is rejected: \code{g} is used after the push, and it holds \code{xs}. The cost is linear.

\section{Checking, and reading off the summary}
\label{sec:inf-summarise}

At each demand, each call of a once closure and each creation of one, the four side conditions are
applied. Then, for each tracked parameter path $\pi_i.q$:
\begin{itemize}
\item $\mathsf{W} \in \mathit{use}(i, q)$ if some demand's operand may hold $\pi_i.q$ (the local holders
were discharged by \Dead; the caller discharges the rest);
\item $\mathsf{A} \in \mathit{use}(i, q)$ if $\pi_i.q \in A(\mathit{result})@q'$ for some $q'$;
\item $\mathsf{E} \in \mathit{use}(i, q)$ if $\pi_i.q$ was stored in an escaped place, passed to a
position whose $\mathit{use}$ contains $\mathsf{E}$, captured by an escaping closure, or passed to a
\code{Js} operation or a \code{foreign} with no asserted summary;
\item $\mathsf{R} \in \mathit{use}(i, q)$ if some instruction dereferences a holder of $\pi_i.q$, or
passes it to a position whose $\mathit{use}$ contains $\mathsf{R}$;
\item $\mathit{use}(i, q) = \unused$ if none of these holds;
\item $\mathit{res}(q') = \{\pi_i.q \mid \pi_i.q \in A(\mathit{result})@q'\} \cup \{\fresh_k \mid
\site(\iota_k) \in A(\mathit{result})@q'\}$, numbering the sites in instruction order, so that two
result paths that may hold one new cell name the same $\fresh_k$;
\item $\mathit{xres}(q')$ and $\mathit{xreq}(i, q)$ from the flag of \S\ref{sec:model-excl};
\item $\mathit{fun}(j)$ from the calls of $\pi_j$ the walk met: $\mathit{calls}$ by counting along
paths (a call inside a tail-call loop, or a pass to a $\mathit{many}$ position, is $\mathit{many}$);
$\mathit{supply}(k)$ from the abstract value of the $k$th argument at each call, its provenance being the
parameter paths and fresh sites in it and its ownership bit owned iff the argument's cells pass \Dead,
\Unescaped\ and \Disjoint\ at that call; $\mathit{need}$ from whether the result's cells reach a store
that promises $\excl$.
\end{itemize}
The four bits are independent; the normal shape of a writer that returns its operand is
$\{\mathsf{R}, \mathsf{A}, \mathsf{W}\}$. The bits refine Aspinall, Hofmann and Kone\v{c}n\'y's three
usage aspects, in which variables ``may only be accessed once with aspect 1'', ``can be used many times
with aspect 2'', and ``can be freely used with aspect 3'' \citep[\S1]{aspinall2008usage}, with ``used
once'' replaced by ``unique at the site'', and with a bottom element, $\unused$, that their system has
no need for.

\section{Recursion: an optimistic fixpoint}
\label{sec:inf-recursion}

A strongly connected component of the call graph is solved together. (beni's module graph is
acyclic, so recursion never crosses a module boundary.) Every member starts from the most
optimistic summary, the bottom of the lattice of \S\ref{sec:inf-lattice}: every parameter path
$\unused$, every result path empty, every container exclusive and no exclusivity required, every
function-typed parameter called at most once, supplied owned cells and needing nothing. The bodies are
analysed under these assumptions, the resulting summaries are joined with the assumptions, and the
process repeats until nothing changes. Every transfer is monotone---a more pessimistic callee summary
creates more holders, aliases and escapes in the caller, never fewer---and the lattice is finite, so
the iteration terminates (\S\ref{sec:inf-lattice} gives the height). Because no group can fail to
converge, a cap on rounds, if an implementation wanted one, would yield the top summary and a limit
report, never a wrong answer.

\emph{The group is then validated}: each body is checked once more against the final summaries, and
every side condition must pass. A summary assignment that survives validation is a consistent typing
of the group, and Lemma~\ref{lem:summary} argues its soundness by induction over executions.

This is the procedure the literature converges on. Aspinall, Hofmann and Kone\v{c}n\'y report a
result of Kone\v{c}n\'y's \citep{konecny2003typing} that ``for every program typable in our system,
there is a typing with best (i.e.\ largest) usage aspects'', which ``can be automatically
reconstructed using an iterative search for a fixed point, starting with the most optimistic typing
of all functions in the program'' \citep[\S6.2]{aspinall2008usage}; we have not been able to check
Kone\v{c}n\'y's original proof, and their system is first-order. Lean~4's borrow inference starts from
the approximation that every parameter is borrowed and marks parameters owned, repeating ``until we
reach a fix point'' \citep[\S5.2]{ullrich2019beans}. Brandon et al.\ observe that each group of their
borrow equations ``defines a monotone function'', justify the existence of a least fixed point by
Knaster--Tarski and termination by finiteness, and note that the equations ``can be viewed as a
Datalog program'' \citep[\S4.4]{brandon2026borrows}. Their argument is for their equations, and ours is
\S\ref{sec:inf-lattice}. beni's write-set analysis solves its own summaries the same way, by Kleene
iteration from the bottom element.

Clarke et al.'s survey explains why ownership-type inference is hard in object-oriented languages:
ownership annotations ``are mostly design-driven'', a trivial assignment always satisfies the
constraints, and qualifier inference therefore ``cannot give a best solution''
\citep[\S5]{clarke2013survey}---a satisfying solution need not be the intended one. Our facts are not
design choices but properties of the program, so the least solution of the monotone constraints
\emph{is} the intended one: it is the most permissive summary every body validates against. The
closest work in that survey is Huang et al.'s Kleene iteration of a monotone transfer function from the
lattice's bottom, maximising a programmer-supplied metric \citep[\S5.2]{clarke2013survey}.

In the Clean tradition, by contrast, an expression whose typing would need a choice among instances
is declared untypable---``we do not try any specific instances but consider the expression
untypable'' \citep[\S8]{barendsen1996uniqueness}---and the type of every binding in a recursive binding
group is non-unique \citep[\S7]{devries2008simplified}, a constraint on the recursively bound names
themselves rather than on their parameters. Our unknowns are facts about flows rather than attributes
on types; a recursive function's own name is a reference to a declaration that captures nothing, so it
holds no cell, and a recursive call is analysed through the group's summary like any other call.

\section{The lattice}
\label{sec:inf-lattice}

Termination and the existence of a least solution need the whole order, with every fact oriented so
that ``up'' is ``more pessimistic''. For one declaration with parameter paths $P$, result paths $Q$
and function-typed parameters $J$, the facts and their orders are:
\begin{center}\small
\begin{tabularx}{\linewidth}{@{}>{\raggedright\arraybackslash}p{0.3\linewidth}>{\raggedright\arraybackslash}p{0.3\linewidth}X@{}}
\toprule
Fact & Values, bottom first & Height \\
\midrule
$\mathit{use}(i, q)$, $\pi_i.q \in P$ & subsets of $\{\mathsf{R}, \mathsf{A}, \mathsf{E}, \mathsf{W}\}$, by inclusion & 4 per path \\
$\mathit{res}(q')$, $q' \in Q$ & subsets of $P \cup \{\fresh_1, \ldots, \fresh_{|Q|}\}$, by inclusion & $|P| + |Q|$ per path \\
$\mathit{xres}(q')$ & $\mathsf{yes} \sqsubset \mathsf{no}$ (a \emph{must} fact, reversed) & 1 per path \\
$\mathit{xreq}(i, q)$ & $\mathsf{no} \sqsubset \mathsf{yes}$ & 1 per path \\
$\mathit{calls}(j)$ & $0 \sqsubset {\le}1 \sqsubset \mathit{many}$ & 2 per $j$ \\
$\mathit{supply}(j, k)$ & provenance sets by inclusion; $\mathrm{owned} \sqsubset \mathrm{borrowed}$ (reversed) & $|P| + 2$ per parameter of $g_j$ \\
$\mathit{need}(j)$ & nothing $\sqsubset$ fresh and exclusive & 1 per $j$ \\
once-ness of a result path & not once $\sqsubset$ may be once & 1 per function-typed result path \\
\bottomrule
\end{tabularx}
\end{center}
The base lattice $L_0$ is the product of these, ordered pointwise, and its height $h_0$ is the sum of
the last column. A \emph{conditional} summary maps each valuation of its guard atoms to a point of
$L_0$, monotonically: a more pessimistic class for a function-typed parameter never yields a more
optimistic summary. With $a$ atoms---facts of the function-typed parameters' classes that the body's
facts depend on---the normal form is a join of at most $2^a$ guarded values, one per conjunction of
atoms, and its lattice has height at most $2^a \cdot h_0$. We cap $a$ (at eight, say): a body whose
facts would depend on more atoms is summarised by its unconditional, most pessimistic value for the
excess atoms, with a limit tag---the same shape as the cap on inferred constraints that beni already
has. The group's lattice is the product over its members, whose height is the sum of theirs.

Every transfer of Table~\ref{tab:transfer} and every summary rule of \S\ref{sec:inf-summarise} is
monotone in this order: each is a union, a projection, a membership test on a set that only grows, or a
must fact that only falls. We have checked this case by case but have not written the case analysis
out; it is the first item of a mechanised proof. Given monotonicity, Knaster--Tarski gives a least
fixpoint, Kleene iteration from bottom reaches it in at most the group's height in rounds, and the
optimistic start finds it. In practice one or two rounds suffice, because summaries move only when a
recursive call carries a new fact around the cycle.

\section{Function values: ownership classes}
\label{sec:inf-classes}

A function-typed parameter $g$ of $f$ has no body to analyse. One option is to assume the worst---every
argument escaped, the result $\top$, called many times, possibly once---which is the position of
Huisinga's static uniqueness analysis for Lean, where ``higher-order functions are always shared''
\citep[\S4.1]{huisinga2023lean}, a limitation that thesis itself calls considerable. The other is to
make $f$'s summary \emph{polymorphic} in $g$'s. beni already does the second for effect bits: a
reference to a top-level declaration instantiates its summary, each class of the scheme receiving a
fresh class at the use, related to the scheme's classes by inclusion edges. That is what lets one
\code{List.map} serve both a pure and a suspending callback. Ownership takes the same place:

\begin{itemize}
\item every function type in a scheme---a parameter, a result, a field, an element, a captured
value---carries an \emph{ownership class}: the $\mathit{use}$, $\mathit{res}$ and $\mathit{xres}$ facts
of the function type's own parameters and result, whether its values may be once, and what they may
hold under $\env$;
\item $f$'s summary is expressed over its function-typed parameters' classes, as the guarded values of
\S\ref{sec:inf-lattice};
\item at a call \code{f ⋯ λ ⋯}, the lambda's actual summary is related to the instantiated class by
\emph{directional edges}, not unification:
\begin{itemize}
\item the lambda may write its $k$th parameter only where $f$ supplies an owned cell:
$\mathsf{W} \in \mathit{use}_\lambda(k) \Rightarrow \mathit{supply}_f(k) = \mathrm{owned}$;
\item a once lambda requires $\mathit{calls} \le 1$;
\item the lambda's result is translated into $f$'s vocabulary by substituting, for each parameter
$\pi_k$ of the lambda named in $\mathit{res}_\lambda$, the provenance $\mathit{supply}_f(k)$, and the
translation flows into $f$'s $\mathit{res}$ wherever $f$ returns or stores $g$'s result;
\item the lambda's captured cells under $\env$ are the cells of the function-valued argument, so they
take part in \Disjoint\ at the call, and they become escaped if $f$'s class puts $\mathsf{E}$ on that
argument's $\env$;
\end{itemize}
\item a violated edge is an error at the call, naming the lambda and the parameter. As with effects,
this is subtyping, not just unification.
\end{itemize}

The translation through $\mathit{supply}$ is what an earlier version of these rules lacked. There,
$\mathit{supply}$ said only ``owned'' or ``borrowed'', so in
\begin{Verbatim}
pick rows = List.foldl rows [] (λr acc → r)
\end{Verbatim}
the lambda's $\mathit{res} = \{\pi_1\}$ named the lambda's own parameter, which nothing at the call of
\code{foldl} could relate to \code{rows}. With provenance, \code{foldl}'s summary
\begin{Verbatim}
foldl : List a [borrowed; [*] aliased], b [consumed+aliased if g consumes 2; aliased otherwise],
        (a, b → b) [calls many; supply 1 borrowed π₁[*]; supply 2 owned π₂ ∪ g.res]
          → b [= π₂, g.res]
\end{Verbatim}
instantiates the lambda's $\{\pi_1\}$ to $\{\pi_1.[\ast]\}$, so \code{pick}'s result aliases the
elements of its argument, and a caller that edits the result while still using \code{rows} is
rejected.

The classes live on solved types, so a function value that reaches $f$ through a record field or a
parameter of a parameter is handled by the same instantiation: the class is where the function type
is. Rank-1 classes suffice. Lorenzen et al.\ note that ``higher-order functions in Rust that pass newly
borrowed values to their callbacks must have higher-rank types \ldots\ while in our system such
functions have rank-1 types that can be fully inferred'' \citep[\S8.3]{lorenzen2024oxidizing}; beni's
arrows are $n$-ary and saturated, so there is no curried spine along which to propagate.

What this does \emph{not} do is infer, inside $f$, that a particular closure only reads a capture and
may therefore be called many times while the capture is consumed later. That fact belongs to the
lambda (its $\env$ cells read, not written) and is decided where the lambda is written; it composes.
In \code{List.map xs (λx → f x ys)} the lambda reads \code{ys}, \code{map}'s $\mathit{calls} =
\mathit{many}$ is fine for a closure that is not once, and the closure does not escape, so \code{ys}
is live only during the \code{map} call and a later \code{List.set ys ⋯} is accepted. Deng et al.\ draw a
related distinction with \emph{use} and \emph{kill} effects on closure types: a closure that accesses
a resource, by reading or writing it, carries a use effect on it, and one that frees it a use and a kill
\citep[Fig.~1, \S2.2]{deng2025capture}. Their use effect corresponds to our ``not consumed'', not to
``read only''. System Capybara draws a different line: it accepts a write between two calls of a reading
closure, and rejects ``only the interfering use''---a write run in \emph{parallel} with the reading
closure \citep[\S2.2]{xu2026capybara}. Capybara is more permissive than we are in sequential code,
since an in-place write must not be observed by a closure that is still live, and it infers read-only
permissions from capture sets while still requiring declared markers on mutable classes and on
mutating and consuming methods \citep[\S6.1]{xu2026capybara}. We infer the distinction because the
lambda is in front of the checker.

\section{Constraints and generic code}
\label{sec:inf-generic}

A method call \code{x.m a} inside a function with a \code{where} clause calls evidence the caller
supplies. There are two cases.
\begin{itemize}
\item \emph{The well-known methods \code{eq} and \code{compare} only read.} An earlier version of these
rules said so ``by contract'' while also saying that a user \code{eq} is summarised like any function;
the two contradicted each other, and a user \code{eq} that pushed onto its argument, reached through
\code{==} on a list of that type, wrote the same cell twice behind the program's back. We close the
case with a language rule: \emph{a type's \code{eq} and \code{compare} must have $\mathit{use} \subseteq
\{\mathsf{R}\}$ on every parameter path}, and one that may write, keep or return a cell of its arguments
is an error at its declaration. Derived \code{eq} and \code{compare} satisfy it by construction, and
the library's list equality and comparison, which call the element's evidence, then only read. The rule
costs a programmer nothing they could not write: an \code{eq} that wants to sort a copy calls
\code{List.sort}, which builds a new list, or \code{List.copy} before an in-place writer. It buys a
guarantee---\code{==} and \code{<} never change a value---on which every generic container's summary
then rests (\S\ref{sec:disc-options}, Change~11 weighs it). The alternative, summarising \code{eq} like
any method and making every \code{where a.eq} function's summary conditional on its evidence, is sound
too; it costs every generic function a guard on its evidence class and makes \code{List.eq}'s asserted
row conditional, for a capability we know no use for.
\item \emph{Any other method}, such as \code{a.update : a, Int → a}, has a summary that is a class on
the constraint's method type, instantiated at each call site from the term the dispatch table
resolves---a top-level function, a derived method, or evidence received as a parameter---exactly as its
effect class is. Inside the generic body the class's variables are what the body's summary is written
over; at the use they are bound.
\end{itemize}

Generic code therefore loses nothing by being generic. What it cannot do is write a value of type
\code{a} in place: there is no demand of type \code{a → a}. A generic container function such as
\code{Dict.update d k (λv → List.push v x)} is the container's business (\S\ref{sec:hard-containers}),
and its callback's $\mathit{supply}$ is a class against which the caller's lambda is checked.

Parametricity supplies result aliasing for free---a function of type \code{a → a} can only return its
argument or something a callback gave it, which is Futhark's observation about \code{id} and the pipe
\citep{henriksen2026aliasing} and the identity typing in reachability types \citep{wei2024polymorphic}.
The checker obtains the same fact from the body (\code{identity x = x} analyses to $\mathit{res} =
\{\pi_1\}$), and from the type for a \code{foreign} over a type variable, such as \code{Debug.log},
\emph{only when no opaque runtime type with a tracked parameter stands between argument and result}.
\code{Task.join : Fiber a → a} returns a value the runtime retains and hands to every joiner, so its
result is $\top$, not ``fresh from nowhere''.

\section{Complexity}
\label{sec:inf-complexity}

Let $n$ be the number of instructions of a declaration, $w$ the number of tracked paths of the widest
type it mentions, $c$ the number of distinct abstract cells in it, $D$ its demands and $H$ the
holders of a cell at a demand.
\begin{itemize}
\item \emph{Width.} With recursive types folded (\S\ref{sec:model-cells}), $w$ is at most the number
of (type, step) edges reachable in the definitions of the types involved, not exponential in the cut:
a tree \code{type Tree = Node String (List Tree)} contributes one list path and one element path. For a
non-recursive type, $w$ is bounded by its fields to depth $k$; a model record with three list fields
has $w = 3$.
\item \emph{Forward and backward walks.} $O(n \cdot w)$ each per round; a self tail call adds an inner
fixpoint of at most $c$ rounds, so a declaration with a loop costs $O(c \cdot n \cdot w)$.
\item \emph{The check.} $O(D \cdot H)$, which is quadratic in the size of one declaration in the worst
case, when a cell has many alias bindings and many demands. In a TEA application the largest
declaration is usually \code{update}, which may be most of a module, so the honest statement is that
the check is quadratic in the size of the largest declaration.
\item \emph{Recursion.} A group runs at most as many rounds as its lattice's height
(\S\ref{sec:inf-lattice}), which grows with the product of result paths and parameter paths through
$\mathit{res}$, and with $2^a$ for $a$ guard atoms; in practice one or two rounds.
\end{itemize}
Interfaces grow by a few words per scheme site, plus the crossing bit per type. Nothing is
whole-program: summaries are read from interfaces exactly as effect summaries are, and the entry
protocol of a TEA program (\S\ref{sec:hard-tea}) is evaluated at the one call that starts the program,
from summaries.

The compiler's design target for checking is above 250\,000 lines per second per core, and the effect
inference, which has the same shape with a narrower per-site record, fits within a 10\,\% increase in
check time. We expect the ownership pass to fit similarly, but have not measured it: the unknowns are
$w$ on wide model records, $H$ inside a large \code{update}, and the number of guard atoms, and the
evaluation (Chapter~\ref{sec:evaluation}) measures all three. If a cap is needed on $w$, a type with
more than a fixed number of tracked paths (say 64) is tracked as a whole and yields $\top$---a limit
report and never a wrong answer. A cap of 64 is not expected to fire on recursive types once they are
folded; it would fire on very wide records.

\section{Modularity and determinism}
\label{sec:inf-modular}

The summary of every public declaration is published in the module interface as a block beside the
effect facts: per scheme site, found by the same encoding of a site as a root and a path that the effect
block uses, the facts of \S\ref{sec:model-summary}, and per function-type site its class. The block
also carries the crossing bit of every type the module declares. The block is \emph{hashed with the
interface}. A body edit that changes a summary therefore changes the hash, and importers are
re-checked---the only correct behaviour, since their side conditions change. The price is interface
churn: with inferred summaries, editing a function's body can change its signature, which is the usual
argument for requiring signatures to be written. In our measurement of beni's inferred \code{where}
clauses, one class of body edit changed the interfaces of 50 of 90 unannotated public declarations of
the core library and of none of the annotated ones (\S\ref{sec:bg-modules}); we expect summaries to
behave similarly, and the evaluation measures it. An optional annotation can freeze a summary
(\S\ref{sec:disc-options}, Change~5).

\emph{Everything the analysis reads about another module is in that module's interface.} This
includes the entry protocol of a TEA program: it is evaluated at the call that starts the program, from
the summaries of \code{init} and \code{update} and the other functions the program record names, and
those summaries carry exactly the facts the entry needs---numbered fresh cells for aliasing between
result paths, and $\mathsf{E}$ for anything a command or subscription keeps
(\S\ref{sec:hard-tea}). If \code{init}'s body changes from building two lists to placing one list at
two fields, its $\mathit{res}$ changes from $\{\fresh_1\}, \{\fresh_2\}$ to $\{\fresh_1\}, \{\fresh_1\}$,
the interface hash changes, and the module that starts the program is re-checked. An earlier version
of these rules computed the entry state from bodies the interface did not carry, and a cached
acceptance could then outlive such an edit. The cache key therefore does not change: the summary block
is a function of the module's source and its imports' interfaces, which the key already covers. The
incrementality tests gain this case---an aliasing edit in \code{init} in one module, the program
started in another---and the determinism test the ownership dump.

beni requires that its output not depend on thread timing (\S\ref{sec:bg-modules}). Every order the
pass uses is the checker's: components in source order, instructions in order, fields and paths by
name, classes by their index in the type store, fresh cells by instruction index. No pointer is hashed
and no thread order is observed. The determinism test gains the ownership dump of
\S\ref{sec:eval-report} beside the dumps it already compares, with cells printed as
\code{site <decl>:<inst>} and \code{param <i>.<path>}, both derived from the source, summaries in
scheme-site order and holders in source order.

\section{What the pass is not}
\label{sec:inf-not}

It is not an optimisation that decides per site whether to write in place: every accepted demand is
in place, so the backend needs no per-site bit (\S\ref{sec:compiler-backend}). Nor is it the
optimiser's analysis. Its verdict is a function of the source and the rule set; inlining and
dead-binding elimination never change it. The optimiser is instead \emph{told} where the demands are
and moves no read across one (A5). This is the condition McCall set for any such guarantee in Swift:
it must be a real guarantee, ``even in debug builds'', and never be defined by what a particular
optimiser can eliminate in a particular release \citep{swift2025explicitcopies}.
